\section{Background}
\label{sec:background}

\subsection{3D Gaussian Splatting and Feed-Forward Reconstruction}

3DGS represents a scene as a set of anisotropic 3D Gaussians, each with a position, a shape, an opacity and a color, and renders a view by projecting the Gaussians onto the image plane and blending them in depth order~\cite{kerbl2023gaussian}.
Kerbl et al.\ fit these parameters to each scene by gradient descent from posed images, after the capture is complete.
Feed-forward models replace this optimization with one pass of a trained network that predicts the Gaussians from the input frames.

AnySplat takes a set of unposed video frames and predicts, in one call, a set of Gaussians together with the camera poses of the frames~\cite{anysplat2025}.
Each call reconstructs $k$ frames and returns a fresh set of Gaussians, so successive calls share no Gaussian identities.
In our measurements its cost grows with $k$: from $k=2$ to $k=16$ one call went from 0.12 to 0.86~s and the process's GPU memory from 4,726 to 9,576~MiB. % G1

StreamSplat reconstructs a dynamic scene online from an uncalibrated monocular video~\cite{streamsplat2026}.
It resizes each frame to $512\times288$ and pairs it with a relative depth map from Depth Anything V2 ViT-L~\cite{depthanythingv2}. % G5
A static encoder, a vision transformer that sees one frame at a time, predicts pixel-aligned Gaussians, 36,864 per frame. % docs/reading_notes.md (256x144 grid); G5 counts 73,728 per pair
A dynamic decoder then attends across a pair of consecutive keyframes and predicts, for every Gaussian of both frames, a velocity and two coefficients that fade its opacity over time, so the pair can be rendered at any time between its two keyframes.
The paper's online algorithm encodes each new keyframe once and keeps the previous keyframe's features for the next pair, while the released inference code encodes every interior keyframe twice.
One update in our server is one new keyframe: a depth map for it, then Gaussians and motion for the pair formed by the previous and the new keyframe, 73,728 Gaussians in all. % G5
The paper reports 49~ms per frame on an NVIDIA A100 without saying whether this includes depth; on our GPU the released code takes 0.19~s per update, 0.049~s of it for depth. % 49 ms, A100, depth unstated: docs/reading_notes.md (paper Table 8), G4; 0.19 and 0.049 s: G5 (G5's "unnamed GPU, excludes depth" disagrees with the notes)

The StreamSplat paper evaluates on DAVIS~\cite{davis2017} with keyframes five frames apart and computes PSNR at $256\times256$ on two sets of frames. % G8
Keyframe PSNR re-renders the two keyframes of a pair, which the model was given, so it measures how faithfully the Gaussians encode the input.
Middle-frame PSNR renders the four frames between the keyframes, which the model never sees, from the predicted motion, so it measures interpolation in time. % G8
The two differ by about 14~dB, both in the paper (37.83 and 23.66~dB) and in our reproduction on DAVIS 2017 val (37.76 and 23.76~dB). % G8
Across the 30 val sequences, our per-sequence mean keyframe PSNR ranges from 33.96 to 42.40~dB and the middle-frame PSNR from 17.09 to 34.61~dB. % G8

\subsection{Sharing One GPU Among Several Users}

When several processes use one NVIDIA GPU, each has its own context, and by default the GPU time-slices among the contexts, so kernels from different processes do not run at the same time~\cite{nvidia_mps}.
This needs no setup, and it is what independent per-user processes get.

MPS routes the work of several client processes through one server process that holds a single GPU context, so kernels from different clients can run concurrently~\cite{nvidia_mps}.
MPS can also cap each client's share of the streaming multiprocessors and of device memory.

MIG splits a supported data-center GPU into instances with their own compute units, memory and memory bandwidth, isolating tenants in hardware~\cite{nvidia_mig}.
The RTX 6000 Ada used in this report does not support MIG, and MIG was disabled on the H200 nodes available to us.

An inference server can also share a GPU at the application level: one process holds one copy of the model weights, collects the pending requests of several users and runs them as one batch~\cite{clockwork2020}.
Batching helps when one request leaves the GPU underused, and each request then waits for its whole batch.

A CUDA graph records a sequence of kernel launches and memory operations once and replays it with a single launch, which removes the CPU cost of launching each kernel~\cite{nvidia_cuda_guide}.
Work captured in a graph cannot wait on the host, so host synchronizations inside the model must be removed before capture. % G13

In a cluster, Kubernetes assigns GPUs to containers through NVIDIA's device plugin, which advertises each GPU to the scheduler as a resource that pods request~\cite{k8s_device_plugin}.
In its default configuration a pod receives whole GPUs, and the plugin's sharing options build on time-slicing and MPS.
